\documentclass{article}
\usepackage[utf8]{inputenc}
\usepackage[T1]{fontenc}
\usepackage{geometry}
\geometry{
    a4paper,
    left=2.5cm,
    right=2.5cm,
    top=2.5cm,
    bottom=2.5cm
}
\usepackage{graphicx}
\usepackage{float}
\usepackage{hyperref}
\usepackage[english]{babel}
\usepackage{amsmath}
\usepackage{amssymb}
\usepackage{array}
\usepackage{booktabs}
\usepackage{caption}
\usepackage{subcaption}
\usepackage{enumitem}
\usepackage{fancyhdr}
\usepackage{titlesec}
\usepackage{url}
\usepackage{textcomp}
\usepackage{xcolor}
\usepackage{listings}
\usepackage{multirow}
\usepackage{longtable}
\usepackage{tabularx}
\usepackage{times}

% Set up page style
\pagestyle{fancy}
\fancyhf{}
\rfoot{\thepage}
\renewcommand{\headrulewidth}{0pt}

% Title formatting
\titleformat{\chapter}[display]{\normalfont\bfseries}{\chaptertitle\thechapter}{0pt}{\Huge}
\titlespacing*{\chapter}{0pt}{50pt}{40pt}

% Document info
\title{Template Document}
\author{Academic Support Agent}
\date{\today}

\begin{document}
\maketitle

\maketitle
\begin{enumerate}[leftmargin=\labelsep]
\item Question summary/restatement can go here.
    \begin{enumerate}
    \item Summary of the first part of the question can go here.

    Down here is then where you can place your solution. Math can be entered using the equation
    environment like this
    \begin{equation}
        \vec{\mathbf{r}} = \vec{\mathbf{r}}_{0} + \vec{\mathbf{v}}_{0}t + \frac{1}{2}\vec{\mathbf{a}}t^{2}
    \end{equation}
    If you then where working in say the $x$-direction and had some numbers % A percent sign allows you to comment.
    %The dollar signs around something in a line of text is for "in-line math"
    \begin{equation}
    \begin{array}{r@{~=~}l}
    x & x_{0} + v_{x0}t + \frac{1}{2}a_{x}t^{2} \\ [2ex]
    & 1.2~\text{m} + (4.0~\text{m/s})(3.0~\text{s}) + \frac{1}{2}(-1.0~\text{m/s}^{2})(3.0~\text{s})^{2}\\ [2ex]
    & \boxed{8.7~\text{m}}
    \end{array}
    \end{equation}
    Here we used an array so that we could have multiple lines in a single equation environment. The
    \verb"@{~=~}" sets things up so that when we type a \& it will put in a space, then and equals sign
    and another space. The r means to take what is on the left of the \& and align it to the right
    and the l then means to left align what follows the \&.

    \item When you get to the next part, just add a \verb"\item" to get the appropriate label. Also,
    if you don't like all the equation numbers, you can use the following to have the equation with
    no number
    \begin{equation*}
    \sum\vec{\mathbf{F}} = m\vec{\mathbf{a}}
    \end{equation*}

    \item For more details on putting math into {\LaTeX} documents you can see 
    \href{https://www.overleaf.com/learn/latex/Mathematical_expressions}{this page on Overleaf}.
    \end{enumerate}

\item We you get to the next problem, just make sure to end the enumerate for the parts of the previous
problem and then add another item.
    \begin{enumerate}
    \item Use a nested enumerate environment to label the parts of the next problem.
    \item For a quick and broad overview of how to create documents in {\LaTeX} see 
    \href{https://www.overleaf.com/learn/latex/Learn_LaTeX_in_30_minutes}{this quick tutorial from Overleaf}.
    \end{enumerate}
\end{enumerate}
\end{document}